\section{Main Results}
\label{sec:results}

We test whether the CECL Day-One disclosure triggered a disciplinary
response.  Three predictions follow: banks whose CECL adjustment revealed
larger latent credit losses should experience relative declines in
uninsured deposit balances after adoption, should face higher funding
costs to replace those balances, and should suffer lower profitability.
The results confirm all three.  Table~\ref{tab:headline_summary} collects
the headline post-adoption effects at high-CECL banks in a single panel
for reference; the remainder of this section develops the individual
estimates.  We anchor on the within-bank composition shift between
uninsured and insured deposits---identified under bank-by-quarter fixed
effects---because it is the estimate that survives every specification
we throw at it in Section~\ref{sec:robustness}.

Two clarifications on scope.  The primary outcome is uninsured
\emph{time} deposits, the deposit category most susceptible to
disciplinary withdrawal at the rollover margin and the category with the
sharpest theoretical prediction \citep{maechler2006dynamic,
egan2017deposit}.  Total uninsured deposits---which include transaction
balances---move in the same direction but with weaker statistical
resolution because they mix disciplinary responses with routine
liquidity flows; we report time deposits as the cleaner test of the
mechanism, not because the effect is confined to that line item.  Second,
the aggregate implication is modest by design: 217 high-CECL banks with
mean assets of \$746~million and a 0.394 percentage-point uninsured
outflow imply roughly \$0.64~billion in reallocated funding, orders of
magnitude smaller than the tens of billions withdrawn from the failed
institutions in March 2023.  Our claim is not that CECL disclosure
substitutes for the 2023 panic in economic magnitude; it is that a
mandatory informational disclosure, absent a panic trigger, moves
uninsured funding measurably and durably at community banks, with
detectable pass-through to funding costs, profitability, and credit
supply.

\subsection{Composition of Time Deposits}
\label{sec:results:time_deposits}

Time deposits are the deposit category most susceptible to disciplinary
withdrawal: certificates are held for investment rather than transaction
purposes, so discipline operates at the rollover margin
\citep{calomiris1991role, iyer2016tale, maechler2006dynamic,
egan2017deposit}.

We stack the insured and uninsured time-deposit series in a single panel
and estimate a triple difference:
\begin{equation*}
\begin{split}
  y_{idt}
  = \alpha_{it}
  + \beta_1 \cdot \mathrm{Uninsured}_d
  + \beta_2 \cdot \mathrm{Post}_{t} \times \mathrm{Uninsured}_d
  + \beta_3 \cdot \mathrm{CECL}_i \times \mathrm{Uninsured}_d \\
  \quad + \beta_4 \cdot \mathrm{CECL}_i \times \mathrm{Post}_{t} \times \mathrm{Uninsured}_d
  + \varepsilon_{idt},
\end{split}
\end{equation*}
where $d \in \{\text{uninsured}, \text{insured}\}$ and $\alpha_{it}$ are
bank $\times$ quarter fixed effects.  These absorb the full vector of
bank-specific time-varying shocks in each period, so identification comes
from the differential movement of uninsured relative to insured deposits
within the same bank and quarter.  The triple-interaction coefficient
$\hat{\beta}_4$ captures the within-bank composition shift.

Table~\ref{tab:small_bank_triple_diff_time_deposits} reports the estimates.
The binary triple interaction (column~2) is $-1.252$ (SE $0.340$,
$p<0.01$): relative to their own pre-adoption composition and relative to
low-CECL banks, high-CECL banks experienced a 125-basis-point widening of
the insured-to-uninsured gap as a share of assets after adoption.  The
continuous analog (column~1, $-0.177$, $p<0.01$) confirms a smooth
dose-response.  Bank $\times$ quarter fixed effects absorb any level
shift in total time-deposit demand at each bank in each quarter---
including any bank-specific effect of the 2023 stress, any bank-specific
change in liability strategy, and any bank-specific shift in
loan-to-deposit demand---so identification comes strictly from the
differential movement of uninsured relative to insured deposits within
the same bank-quarter cell.

Two clarifications on what this design does and does not deliver.
Because both arms of the differential are deposit categories, the
comparison is between two potentially endogenous outcomes rather than
between the treated outcome and an unaffected control liability
\citep{covitz2004market}.  In particular, we show in
Section~\ref{sec:robustness:destination} that the insured
increase reflects active bank substitution into insured brokered and
reciprocal deposits, not only spontaneous depositor reallocation; the
triple-difference should therefore be read as the joint depositor+bank
response at the insured/uninsured margin rather than as a
depositor-only quantity.  What the design does deliver is a within-bank
composition contrast that is invariant to any bank-level determinant of
total funding, including the balance-sheet characteristics implicated in
the 2023 stress---which is the property the identification rests on and
which the level DiD does not provide.

The level DiD in Table~\ref{tab:small_bank_did_time_deposits} is
consistent with the composition result: high-CECL banks lost 0.394
percentage points of uninsured time deposits ($p<0.05$, roughly 7 percent
of the sample mean) and gained 0.712 percentage points of insured time
deposits ($p<0.05$).  Continuous specifications match ($-0.072$ and
$+0.070$).  We anchor on the composition contrast because the level
uninsured decline attenuates in Section~\ref{sec:robustness:stress} once
deposit-composition characteristics receive their own time paths---a
bad-control problem that does not affect the within-bank composition
identification.

Figure~\ref{fig:es_time_deps_cecl} plots event-study coefficients for a
ten-quarter window on either side of CECL adoption ($k=-1$ omitted).
Pre-adoption coefficients for uninsured time deposits are clustered near
zero across all ten lead quarters, consistent with parallel trends.  The
post-adoption break emerges within the first adoption quarter and
persists through the end of the sample, while the insured series moves in
the opposite direction with comparable timing.  Extended-window pre-trend
tests are reported in Section~\ref{sec:robustness:pretrends}.

\subsection{Financial Consequences}
\label{sec:results:financial}

Table~\ref{tab:small_bank_did_financial_outcomes} reports DiD estimates
for interest expense, ROA, and NIM.  High-CECL banks paid $0.021$
percentage points more per quarter in interest expense per dollar of
assets ($p<0.01$)---approximately 8--9 basis points annualized, or 9
percent of the sample quarterly mean.  ROA fell by $0.038$ percentage
points per quarter ($p<0.01$)---an annualized 15-basis-point drag---and
NIM contracted by $0.024$ percentage points per quarter ($p<0.05$).  The
joint decline in ROA and NIM is consistent with higher liability costs
not being passed through to borrowers, in line with community banks
facing competitive local loan markets \citep{hannan1988bank}.  Having
lost lower-cost uninsured balances, high-CECL banks replaced them with
insured deposits at rates sufficient to attract depositors who would not
otherwise bear uninsured credit risk---the price channel of depositor
discipline \citep{nier2006market, billett1998cost}.  Continuous estimates
scale proportionally.  Figure~\ref{fig:es_financial_outcomes_cecl} plots
event-study coefficients; pre-adoption leads are flat and the
post-adoption divergence persists through the end of the sample.
